% ASSERT_CONVENTION: natural_units=natural, metric_signature=mostly_minus, fourier_convention=physics, coupling_convention=alpha_s, generator_normalization=T(fund)=1/2, covariant_derivative_sign=D_mu=partial_mu-igA
%
% Lecture 11: The Koide Relation and the Dualised Standard Model
% Part of: The Dualised Standard Model -- Part III Lecture Notes
%
\documentclass[11pt,a4paper]{article}
\usepackage{amsmath,amssymb,geometry,booktabs,xcolor,hyperref}
\usepackage{amsthm}
\geometry{margin=2.5cm}
\newcommand{\Nc}{N_{\!c}}
\newcommand{\Nf}{N_{\!f}}
\newcommand{\SU}[1]{\mathrm{SU}(#1)}
\newcommand{\UU}[1]{\mathrm{U}(#1)}
\newcommand{\fund}{\square}
\newcommand{\antifund}{\overline{\square}}
\newcommand{\adj}{\mathrm{adj}}
\newcommand{\antisym}{\wedge^{\!2}}
\newcommand{\Tr}{\mathrm{Tr}}
\usepackage{tcolorbox}
\tcbuselibrary{skins,breakable}
\newtcolorbox{warningbox}[1][]{enhanced,colback=red!5!white,colframe=red!70!black,fonttitle=\bfseries,title={Important},#1}
\newtcolorbox{previewbox}[1][]{enhanced,colback=green!5!white,colframe=green!50!black,fonttitle=\bfseries,title={Preview},#1}
\theoremstyle{definition}\newtheorem{exercise}{Exercise}[section]

\title{\textbf{Lecture 11: The Koide Relation and the Dualised
Standard Model}}
\author{The Dualised Standard Model --- Part III Lecture Notes}
\date{Lent Term 2027}

\begin{document}
\maketitle

\tableofcontents
\newpage

%% ========================================================================
\section{Introduction and motivation}
\label{sec:intro}
%% ========================================================================

In Lectures~7--8, we constructed the dualised Standard Model: a
magnetic dual of SU(3) QCD with $\Nf = 6$ flavours, where 18~Higgs
doublets emerge from the mes-Higgs field $\Phi_H$ and the fermion
mass hierarchy arises from hierarchical VEVs among these doublets.
The Yukawa coupling $y$ is \emph{universal} (a prediction of the
duality), and the effective Yukawa matrix is
\begin{equation}
  y_{ij}^{\text{eff}} = y\,\frac{\langle H_{ij}\rangle}{v}\,,
  \label{eq:yeff}
\end{equation}
where $H_{ij}$ are the Higgs doublet components and $v$ is the
electroweak VEV.

In this lecture we ask: \emph{does the dualised SM predict or
constrain the fermion mass eigenvalues?}  We establish one solid
result and explore two conjectures:
\begin{enumerate}
\item[\textbf{(R)}] The Koide relation $Q = 3/2$ is algebraically
  equivalent to a \textbf{VEV condition} on the Higgs
  doublet VEV ratios---a reformulation that introduces no new
  parameters and follows from the universal Yukawa structure of
  Lecture~8.
\item[\textbf{(C1)}] The \emph{numerical value} $Q = 3/2$ coincides
  with the SU(5) Casimir ratio
  $C_2(\mathbf{10})/C_2(\mathbf{5}) = 2(N-2)/(N-1)$, which
  uniquely selects $N=5$.  Whether the Higgs potential enforces this
  value is \emph{conjectural}.
\end{enumerate}
The Koide angle $\delta = 2/9$ remains an empirically determined
parameter; its origin is an open problem (see~\S\ref{sec:angle}).

%% ========================================================================
\section{The Koide relation as a VEV condition}
\label{sec:koide-vev}
%% ========================================================================

This section contains the lecture's motivating reformulation.  It requires only
the universal Yukawa coupling of Lecture~8 and standard algebra.
The observation that the Koide condition becomes a VEV constraint
in universal-Yukawa models was first made in the yukawaon
programme~\cite{Koide1983,HabaKoide2008,Ma2007}; we present it
here in the specific context of the dualised SM.

\subsection{Setup: from VEVs to masses}

From Lecture~8, the effective Yukawa coupling for the $k$-th
generation in a given charge sector is
\begin{equation}
  y_k^{\text{eff}} = y\,\frac{\langle H_k\rangle}{v}
  = y\,\frac{\mu_{0k}^2}{M_k^2}\,,
  \label{eq:yk}
\end{equation}
where $\mu_{0k}^2$ is the mixing mass between the SM Higgs $H_0$
and the $k$-th dormant doublet, and $M_k$ is the diagonal mass
of the $k$-th doublet.  The fermion mass is then
\begin{equation}
  m_k = y_k^{\text{eff}}\,\frac{v}{\sqrt{2}}
  = y\,\frac{v}{\sqrt{2}}\,\frac{\mu_{0k}^2}{M_k^2}\,.
  \label{eq:mk}
\end{equation}

Define the \emph{mixing ratio}:
\begin{equation}
  x_k \equiv \frac{\mu_{0k}}{M_k}\,.
  \label{eq:xk}
\end{equation}
Then $m_k = m_0\,x_k^2$ where $m_0 = y\,v/\sqrt{2}$ is the
universal scale, and
\begin{equation}
  \sqrt{m_k} = \sqrt{m_0}\,x_k\,.
  \label{eq:sqrtm}
\end{equation}

\subsection{The Koide condition on mixing ratios}

The Koide parameter for three fermion masses $m_1, m_2, m_3$ is
\begin{equation}
  Q = \frac{(\sqrt{m_1}+\sqrt{m_2}+\sqrt{m_3})^2}
  {m_1+m_2+m_3}\,.
  \label{eq:Q}
\end{equation}
Substituting~\eqref{eq:sqrtm}:
\begin{equation}
  \boxed{Q = \frac{(x_1+x_2+x_3)^2}{x_1^2+x_2^2+x_3^2}\,.}
  \label{eq:Qx}
\end{equation}
The universal scale $m_0$ cancels.  The Koide condition
$Q = 3/2$ is a \textbf{pure constraint on the mixing ratios}
$x_k = \mu_{0k}/M_k$.  This is the key rewriting of this
lecture: in any model with a universal Yukawa coupling and
multiple Higgs doublets, the Koide relation reduces to a
condition on the \emph{Higgs sector} alone.  The cancellation
of $m_0$ is guaranteed by the scale invariance of $Q$
(homogeneous of degree zero), not by any special property
of the dualised SM.

\begin{warningbox}
The Koide condition is independent of the universal Yukawa
coupling $y$ and the Higgs VEV $v$.  It constrains only the
\emph{ratios} $\mu_{0k}/M_k$ of the Higgs sector parameters.
This reformulation is exact.  What follows in Sections~3--4
is conjectural.
\end{warningbox}

\subsection{Decomposition into common and traceless parts}

Write $x_k = z_0 + z_k$ with
\begin{equation}
  z_0 = \frac{x_1+x_2+x_3}{3}\,,\qquad
  z_1+z_2+z_3 = 0\,.
  \label{eq:z0zk}
\end{equation}
Then:
\begin{equation}
  Q = \frac{9\,z_0^2}{3\,z_0^2 + \sum_k z_k^2}\,.
  \label{eq:Qz}
\end{equation}
The condition $Q = 3/2$ requires
\begin{equation}
  z_0^2 = \frac{1}{3}\sum_{k=1}^3 z_k^2\,.
  \label{eq:koide-cond}
\end{equation}
This is the \textbf{trace condition}: the squared common piece
equals one-third of the sum of squared traceless pieces.
It is identical to Koide's original 1983 constraint
$\Tr[U^2] = \Tr[V^2]$~\cite{Koide1983}.

\subsection{What the potential must do}

The effective potential for the dormant Higgs doublets
(Lecture~8) is
\begin{equation}
  V_{\text{eff}} = \sum_k M_k^2\,|H_k|^2
  + \sum_k \mu_{0k}^2\,(H_0^\dagger H_k + \text{h.c.})
  + \text{quartic terms}\,.
  \label{eq:Veff}
\end{equation}
Minimising with respect to $\langle H_k\rangle$ gives
$x_k = \mu_{0k}/M_k$ as before.  The Koide condition~\eqref{eq:koide-cond}
is then a \textbf{constraint on the potential minimum}:
the quartic terms must be structured so that the minimum
satisfies $z_0^2 = \frac{1}{3}\sum z_k^2$.

We do \emph{not} derive that the potential has this property.
The question of whether the emergent quasi-SUSY of the magnetic
theory constrains the quartic couplings sufficiently is the
central open problem of this programme.

%% ========================================================================
\section{Conjecture C1: $Q = 3/2$ and the SU(5) Casimir ratio}
\label{sec:casimir}
%% ========================================================================

\subsection{The numerical coincidence}

For $\SU{N}$, the quadratic Casimir operators of the fundamental
representation~($\fund$) and the antisymmetric two-tensor
($\antisym\fund$) satisfy:
\begin{equation}
  \frac{C_2(\antisym\fund)}{C_2(\fund)}
  = \frac{2(N-2)}{N-1}\,.
  \label{eq:casimir-ratio}
\end{equation}
This follows from $C_2(\wedge^p\,N) = p(N-p)(N+1)/(2N)$.
Setting the ratio to $3/2$:
\begin{equation}
  \frac{2(N-2)}{N-1} = \frac{3}{2}
  \qquad\Longrightarrow\qquad N = 5\,.
  \label{eq:Q-N5}
\end{equation}
This is the unique positive integer solution.

\begin{center}\small
\begin{tabular}{@{}ccc@{}}
\toprule
$N$ & $C_2(\wedge^2 N)/C_2(N)$ & \\
\midrule
$3$ & $1$ & (degenerate) \\
$4$ & $4/3$ & \\
$\mathbf{5}$ & $\mathbf{3/2}$ & $= Q_{\text{Koide}}$ \\
$6$ & $8/5$ & \\
$\infty$ & $2$ & \\
\bottomrule
\end{tabular}
\end{center}

\subsection{Why $N = 5$?}

In the dualised SM of Lectures~7--8, the flavour group is
$\SU{6}$ (from $\Nf = 6$).  The value $N = 5$ does not
correspond to $\SU{6}$ but to $\SU{5} \subset \SU{6}$.
The five lightest quarks $(u,c,d,s,b)$ form the fundamental
of an SU(5) flavour subgroup, with the top quark as the sixth
flavour playing a special role.

This SU(5) flavour structure is precisely the \emph{sBootstrap}
framework~\cite{Rivero2009,Rivero2024}, where the top is called
the ``elephant'' and the five light quarks are ``turtles.''
The decomposition
\begin{equation}
  \SU{6}_{\text{flavour}} \supset \SU{5}_{\text{turtle}}
  \times \UU{1}_{\text{top}}
\end{equation}
is natural from the mass hierarchy alone ($m_t \gg m_b$), but
the duality does not \emph{derive} it: nothing in the CS framework
selects SU(5) over SU(6) as the relevant Casimir group.

\begin{warningbox}
The coincidence $C_2(\mathbf{10})/C_2(\mathbf{5}) = 3/2
= Q_{\text{Koide}}$ is a \emph{mathematical fact} about SU(5).
Whether the Higgs potential minimum is determined by this
Casimir ratio is a \emph{conjecture without proof}.  No
derivation exists connecting the one-loop potential of the
18-doublet sector to the eigenvalue condition~\eqref{eq:koide-cond}.
\end{warningbox}

For this coincidence to be physically meaningful, there must exist
a dynamical channel through which $C_2$ ratios in an $\SU{5}$
flavour group influence the Higgs potential minimum.  Candidate
mechanisms include: one-loop corrections proportional to $C_2$
entering the quartic couplings, renormalisation-group fixed-point
conditions, or representation-theoretic constraints on the
potential structure.  None of these has been investigated.

\subsection{Dynkin index integers}

If one takes the SU(5) Casimir connection seriously, the
Dynkin index ratios provide further numerical coincidences:
\begin{align}
  \frac{T(\mathbf{15})}{T(\mathbf{5})} &= N+2 = 7\,,
  \label{eq:dyn-7} \\
  \frac{T(\mathbf{24})}{T(\mathbf{5})} &= 2N = 10\,,
  \label{eq:dyn-10} \\
  \frac{T(\mathbf{10})}{T(\mathbf{5})} &= N-2 = 3 = \Nc\,.
  \label{eq:dyn-3}
\end{align}
The integers $7$ and $10$ appear in empirical ``$Z$-ladder''
mass relations, and $N-2 = 3$ is the number of QCD colours.
These coincidences are suggestive but do not constitute a
derivation.

%% ========================================================================
\section{The Koide angle $\delta$: an open problem}
\label{sec:angle}
%% ========================================================================

The Koide condition~\eqref{eq:koide-cond} constrains the sum
$\sum z_k^2$ relative to $z_0^2$, but does not fix the
\emph{individual} $z_k$ values.  The mass spectrum is
parametrised by a single angle $\delta$ in the $Z_3$
representation:
\begin{equation}
  x_k = z_0\!\left(1+\sqrt{2}\,
  \cos\!\left(\frac{2\pi k}{3}+\delta\right)\right),
  \qquad k = 1,2,3\,.
  \label{eq:z3}
\end{equation}
The value $\delta = 2/9$~rad reproduces the charged lepton
masses to $10^{-5}$ precision~\cite{Brannen2006}.

Neither the dualised SM nor the sBootstrap currently derives
this value.  Within the multi-Higgs framework, $\delta$ is
determined by the relative magnitudes of the mixing masses
$\mu_{0k}$ and the diagonal masses $M_k$.  These originate
from Planck-scale operators in the electric theory
(Lecture~8), and their structure is not sufficiently
constrained by the duality to fix $\delta$.

Deriving $\delta$ from a dynamical principle---whether from
modular flavour symmetries, conformal fixed-point dynamics,
or some other mechanism---remains the central open problem
of the Koide programme.

%% ========================================================================
\section{What $Q = 3/2$ determines and what it leaves free}
\label{sec:spectrum}
%% ========================================================================

The Koide condition $Q = 3/2$ is one constraint on three
eigenvalues, leaving a two-parameter family: the overall
scale $m_0$ and the angle $\delta$.  The scale $m_0$ sets
the heaviest mass; $\delta$ sets the ratios.

With the empirical $\delta = 2/9$, the Koide formula predicts
\begin{align}
  \frac{m_e}{m_\mu} &= 0.004836\,,
  & \text{(measured: } 0.004836\text{)} \\
  \frac{m_\mu}{m_\tau} &= 0.05946\,,
  & \text{(measured: } 0.05946\text{)}
\end{align}
Both agree to better than $0.002\%$.  But this precision is
the content of the original Koide formula~\cite{Koide1983}; it
is not a prediction of the dualised SM, which so far provides
$Q = 3/2$ as a VEV condition but does not derive $\delta$.

%% ========================================================================
\section{The chain angle and colour}
\label{sec:chain}
%% ========================================================================

The Koide chain~\cite{Rivero2013,Zenczykowski2012}
extends the lepton relation to quarks:
$(0,s,c) \to (-s,c,b) \to (c,b,t)$, with each link
satisfying $Q \approx 3/2$ and the chain angle
$\delta_{scb} \approx 3\,\delta_\ell$.

One interpretation of the factor~$3$: the quark sector carries
SU(3)$_c$ colour, which multiplies the vacuum angle:
\begin{equation}
  \delta_{\text{chain}} = \Nc\,\delta_\ell
  = 3\times\frac{2}{9} = \frac{2}{3}\;\text{rad}
  = 38.2^\circ\,.
  \label{eq:chain-angle}
\end{equation}
The measured value $\delta_{scb} \approx 37.2^\circ$ agrees
to $2.7\%$.  No calculation within the dualised SM or the
sBootstrap currently produces the factor $\Nc$ from dynamics.
Whether the $2.7\%$ agreement is evidence for the colour
interpretation or a numerical coincidence cannot be decided
without a derivation of the vacuum structure in the coloured
sector.

%% ========================================================================
\section{The inverse Koide from the seesaw}
\label{sec:inverse}
%% ========================================================================

\subsection{$Q$ is preserved under inversion}

The Seiberg seesaw (Lecture~3) gives the meson VEV
$\langle M_{kk}\rangle = \Lambda^2/m_k$.  Since $Q$ is
homogeneous of degree zero:
\begin{equation}
  Q\!\left(\frac{\Lambda^2}{m_1},\frac{\Lambda^2}{m_2},
  \frac{\Lambda^2}{m_3}\right)
  = Q\!\left(\frac{1}{m_1},\frac{1}{m_2},\frac{1}{m_3}\right).
  \label{eq:Q-inverse}
\end{equation}
This is a theorem, not a conjecture.

\subsection{Empirical status}

With PDG~2024 pole masses for leptons and $\overline{\text{MS}}$
masses at $\mu = 2$~GeV for quarks~\cite{AHS2025}:
\begin{equation}
  Q(m_e,m_\mu,m_\tau) = 1.50001\,,\qquad
  Q(1/m_d,1/m_s,1/m_b) = 1.505 \pm 0.010\,.
\end{equation}
The uncertainty on $Q(1/m)$ is dominated by $m_d = 4.70 \pm 0.20$~MeV
($4.3\%$).  The deviation from $3/2$ is $0.005 \pm 0.010$, i.e.\
within $1\sigma$: the agreement cannot currently be distinguished
from numerical coincidence at this level of quark-mass precision.

Note that Q-preservation under inversion is a mathematical identity
holding for any mass triplet; the empirical near-agreement
$Q(1/d,1/s,1/b) \approx Q(e,\mu,\tau) \approx 3/2$ is a separate
observation requiring an independent dynamical explanation.

The \emph{individual mass ratios} do not match through a simple
seesaw $m_\ell = c/m_q$ (see~\S\ref{sec:caveats}).

\subsection{The hierarchy problem of the seesaw}
\label{sec:caveats}

The seesaw $m_\ell \propto 1/m_q$ gives an \emph{inverted}
hierarchy: the lightest quark maps to the heaviest lepton.
With the standard generation pairing $(d \leftrightarrow e,\;
s \leftrightarrow \mu,\; b \leftrightarrow \tau)$:
$m_e/m_\mu = m_s/m_d = 19.9$, while the actual ratio is
$0.0048$---a factor of $4000$ discrepancy.

With the reversed pairing $(d \leftrightarrow \tau,\;
s \leftrightarrow \mu,\; b \leftrightarrow e)$: $m_\tau$
is recovered by construction, $m_\mu$ is $15\%$ off, and
$m_e$ is a factor~$4$ wrong.

The seesaw preserves $Q$ (that is exact) but does \emph{not}
reproduce the individual lepton masses.  The quark--lepton
mass mapping is more complex than a simple inversion.

%% ========================================================================
\section{Connection to the sBootstrap}
\label{sec:sbootstrap}
%% ========================================================================

The sBootstrap framework~\cite{Rivero2009,Rivero2024} identifies
sfermions with known mesons and diquarks via $\SU{5}$ flavour
symmetry.  The closure conditions~\cite{Rivero2024}
\begin{equation}
  2N = k_u\,k_d\,,\quad
  2N = \frac{k_d(k_d+1)}{2}\,,\quad
  4N = k_u^2+k_d^2-1
\end{equation}
uniquely give $N = 3$ generations with $k_u = 2$, $k_d = 3$.

In the dualised SM, the ISS rank condition~\cite{ISS2006}
(see also Lecture~6)
provides a \emph{dynamical} derivation of the same split.

\subsection{From $\Nf = 6$ to $\Nf = 5$: the top quark as a trigger}

Start with the full theory: $\SU{3}_c$ with $\Nf = 6$ flavours.
The Seiberg dual is $\SU{\Nf-\Nc} = \SU{3}$ with 6~flavours and
a $6\times 6$ meson matrix $M_{ij} = \tilde Q_i Q_j$ containing
$36$~mesonic composites.  The ISS rank condition at $\Nf = 6$ gives
$\text{rank}(\tilde\varphi\varphi) \leq \Nf - \Nc = 3$,
producing a \emph{symmetric} split: $k_u = k_d = 3$.

Now give one flavour (the top) a large mass $m_t$.  In the
magnetic theory, this corresponds to a VEV for the meson
component $M_{tt}$, which Higgses the dual gauge group:
\begin{equation}
  \SU{3}_{\text{dual}} \;\xrightarrow{\;\langle M_{tt}\rangle \neq 0\;}
  \;\SU{2}_{\text{dual}}\,.
  \label{eq:higgsing}
\end{equation}
Below the scale $m_t$, the effective theory has $\Nf = 5$ active
flavours, the dual gauge group is $\SU{2}$, and the meson matrix
reduces to $5 \times 5 = 25$ components.  The ISS rank condition
on the $\SU{2}$ theory now gives
$\text{rank}(\tilde\varphi\varphi) \leq \Nf - \Nc = 2$:
\begin{equation}
  k_u = \Nf - \Nc = 2 \quad(\text{massless})\,,\qquad
  k_d = \Nc = 3 \quad(\text{massive})\,.
  \label{eq:iss-split}
\end{equation}
The \textbf{asymmetric $2+3$ split is driven by the top mass}:
it is the large $m_t$ that Higgses $\SU{3} \to \SU{2}$ and
breaks the $3+3$ symmetry of the $\Nf = 6$ theory.

\subsection{The meson decomposition}

The $5\times 5$ meson matrix transforms as
$\mathbf{5}\otimes\overline{\mathbf{5}} = \mathbf{24}\oplus\mathbf{1}$
under $\SU{5}$ flavour.  The ISS $2+3$ split decomposes
$\SU{5} \to \SU{2}_u \times \SU{3}_d \times \UU{1}$, and
the 24~adjoint mesons organise by electric charge:
\begin{center}\small
\begin{tabular}{@{}llcc@{}}
\toprule
Block & $\SU{2}\times\SU{3}$ rep & Charge & Count \\
\midrule
$\bar u\,u$\;(traceless) & $(\mathbf{3},\mathbf{1})$ & $0$ & 3 \\
$\bar d\,d$\;(traceless) & $(\mathbf{1},\mathbf{8})$ & $0$ & 8 \\
$\UU{1}$ generator & $(\mathbf{1},\mathbf{1})$ & $0$ & 1 \\
$\bar d\,u$ & $(\mathbf{2},\overline{\mathbf{3}})$ & $+1$ & 6 \\
$\bar u\,d$ & $(\overline{\mathbf{2}},\mathbf{3})$ & $-1$ & 6 \\
\bottomrule
\end{tabular}
\end{center}
giving $\mathbf{24} \to 12_0 + 6_{+1} + 6_{-1}$, plus the
flavour-singlet $\eta'$-like trace for a total of~25.

The sBootstrap~\cite{Rivero2009,Rivero2024} identifies these 25
composites with the scalar partners of known fermions: the charged
mesons ($\bar d_i u_j$ and $\bar u_i d_j$) map to charged
sfermions, while the neutral composites map to sneutrinos and
neutral Higgs-like states.  The closure conditions then
count:
\begin{equation}
  2N = k_u\,k_d = 6\,,\quad
  2N = \tfrac{k_d(k_d+1)}{2} = 6\,,\quad
  4N = k_u^2+k_d^2-1 = 12
\end{equation}
all giving $N = 3$ generations \emph{uniquely}.

\subsection{Koide from the meson mass formula}

The sBootstrap mass formula $m_k = (z_0+z_k)^2\,K_\Omega$
is algebraically equivalent to the Koide condition~\eqref{eq:koide-cond},
as shown in~\cite{Rivero2024}.  The dualised SM provides the
\emph{physical context} for this formula: $z_0+z_k$ is the
mixing ratio $x_k$ of the Higgs sector, and $K_\Omega$ is
the universal scale $m_0$.

%% ========================================================================
\section{Summary and honest assessment}
\label{sec:summary}
%% ========================================================================

\begin{center}
\renewcommand{\arraystretch}{1.3}
\begin{tabular}{lll}
\toprule
\textbf{Result} & \textbf{Origin} & \textbf{Status} \\
\midrule
Koide $=$ VEV condition & Universal Yukawa $+$ algebra
  & \textbf{Reformulated} \\
$k_u = 2$, $k_d = 3$ & ISS rank $+$ top Higgsing
  & \textbf{Derived} \\
$Q(1/d,1/s,1/b) \approx 3/2$ & Seesaw preserves $Q$
  & \textbf{Theorem} \\
$N_{\text{gen}} = 3$ & Closure $=$ duality constraint
  & \textbf{Two derivations} \\
\midrule
$Q = C_2(\mathbf{10})/C_2(\mathbf{5})$ & SU(5) Casimir
  & Observation \\
$\delta_{scb} = \Nc\,\delta_\ell$ & Empirical (factor 3)
  & $2.7\%$ match \\
$\delta = 2/9$ & --- & Open problem \\
\midrule
VEV condition mechanism & Quartic structure
  & \emph{Not derived} \\
Quark--lepton mapping & Seesaw details
  & \emph{Fails quantitatively} \\
$Z$-ladder formula & ---
  & \emph{Not derived} \\
\bottomrule
\end{tabular}
\end{center}

\textbf{Scale convention.}  The VEV condition is imposed at
the scale where the dormant Higgs doublets are integrated out
(near the duality scale $\Lambda \sim$~TeV of Lecture~7).
Below that scale, RG running introduces corrections that spoil
$Q = 3/2$ for quark masses at low energies.  The charged-lepton
Koide relation holds for \emph{pole masses}; the quark
chain uses $\overline{\text{MS}}$ running masses at mixed scales.
This difference is a known subtlety~\cite{Sumino2009} and is
unresolved in the present framework.

\textbf{What is established.}  The universal Yukawa coupling of
the dualised SM translates the Koide relation into a VEV
condition on the multi-Higgs sector.  The ISS rank
condition derives the turtle/elephant decomposition.  The seesaw
preserves $Q$ as a mathematical identity.

\textbf{What is observed but not derived.}  The numerical
coincidence $Q = C_2(\mathbf{10})/C_2(\mathbf{5}) = 3/2$ is an
exact group-theoretic fact that matches the Koide value.
It has no known dynamical origin within the dualised SM.
The angle $\delta = 2/9$ has no known derivation at all.

\textbf{What remains open.}  The quartic coupling structure
of the 18-Higgs potential must be shown to have a minimum at
the Koide point.  The quark--lepton mass mapping through the
seesaw preserves $Q$ but not the individual mass ratios.

%% ========================================================================
\section{Exercises}
\label{sec:exercises}
%% ========================================================================

\begin{exercise}
Verify that $C_2(\wedge^2\mathbf{N})/C_2(\mathbf{N})
= 2(N-2)/(N-1)$ using the general formula
$C_2(\wedge^p N) = p(N-p)(N+1)/(2N)$.  For which values
of $N$ is this ratio an integer?
\end{exercise}

\begin{exercise}
Compute $Q$ for the three charged lepton pole masses
$(m_e = 0.511\;\text{MeV},\; m_\mu = 105.658\;\text{MeV},\;
m_\tau = 1776.93\;\text{MeV})$.  Compare with
$Q(1/m_d,\,1/m_s,\,1/m_b)$ using the PDG values
$m_d = 4.70$, $m_s = 93.5$, $m_b = 4183\;\text{MeV}$.
\end{exercise}

\begin{exercise}
In the $Z_3$ parametrisation~\eqref{eq:z3}, show that
$Q = 3/2$ is satisfied for \emph{any} value of $\delta$.
Then verify that $\delta = 2/9$ reproduces the lepton mass
ratios $m_e/m_\mu$ and $m_\mu/m_\tau$ to better than
$0.01\%$.
\end{exercise}

\begin{exercise}
The sBootstrap closure conditions give $k_u = 2$, $k_d = 3$
uniquely.  The ISS rank condition gives
$k_u = \Nf - \Nc$, $k_d = \Nc$ for SQCD with $\Nf$ flavours
and $\Nc$ colours.  For which values of $(\Nc, \Nf)$ do the
two derivations agree?
\end{exercise}

\begin{exercise}[Challenge]
Write the most general quartic potential for three complex
scalar doublets $H_1, H_2, H_3$ with an $S_3$ permutation
symmetry softly broken by the mixing masses $\mu_{0k}$.
Does the minimum satisfy~\eqref{eq:koide-cond} for any
choice of quartic couplings?  What additional conditions
(if any) are needed?
\end{exercise}

%% ========================================================================
%% BIBLIOGRAPHY
%% ========================================================================
\begin{thebibliography}{20}

\bibitem{Sannino2024}
G.~Cacciapaglia and F.~Sannino,
``Charting Standard Model Duality and its Signatures,''
arXiv:2407.17281 [hep-ph] (2024).

\bibitem{Rivero2009}
A.~Rivero,
``Unbroken supersymmetry without new particles,''
arXiv:0910.4793 [hep-ph] (2009).

\bibitem{Rivero2024}
A.~Rivero,
``An interpretation of scalars in SO(32),''
Eur.\ Phys.\ J.\ C \textbf{84}, 1058 (2024),
arXiv:2407.05397 [hep-ph].

\bibitem{Koide1983}
Y.~Koide,
``New view of quark and lepton mass hierarchy,''
Phys.\ Rev.\ D \textbf{28}, 252 (1983).

\bibitem{Brannen2006}
C.~Brannen,
``The lepton masses,''
\url{https://brannenworks.com/MASSES2.pdf} (2006).

\bibitem{AHS2025}
S.~Antusch, K.~Hinze and S.~Saad,
``Updated Running Quark and Lepton Parameters at Various Scales,''
arXiv:2510.01312 [hep-ph] (2025).

\bibitem{Seiberg1995}
N.~Seiberg,
``Electric--magnetic duality in supersymmetric non-abelian gauge
theories,''
Nucl.\ Phys.\ B \textbf{435}, 129 (1995).

\bibitem{Sumino2009}
Y.~Sumino,
``Family gauge symmetry as an origin of Koide's mass formula,''
JHEP \textbf{0905}, 075 (2009).

\bibitem{HabaKoide2008}
N.~Haba and Y.~Koide,
``New origin of a bilinear mass matrix form,''
Phys.\ Lett.\ B \textbf{659}, 260 (2008),
arXiv:0708.3913 [hep-ph].

\bibitem{Ma2007}
E.~Ma,
``Lepton family symmetry and possible application to the Koide mass formula,''
Phys.\ Lett.\ B \textbf{649}, 287 (2007),
arXiv:hep-ph/0612022.

\bibitem{ISS2006}
K.~Intriligator, N.~Seiberg and D.~Shih,
``Dynamical SUSY breaking in meta-stable vacua,''
JHEP \textbf{0604}, 021 (2006),
arXiv:hep-th/0602239.

\bibitem{Rivero2013}
A.~Rivero,
``A new Koide tuple: strange-charm-bottom,''
arXiv:1111.7232 [hep-ph] (2011).

\bibitem{Zenczykowski2012}
P.~Zenczykowski,
``Remark on Koide's $Z_3$-symmetric parametrization of
quark masses,''
Phys.\ Rev.\ D \textbf{86}, 117303 (2012),
arXiv:1210.4125 [hep-ph].

\end{thebibliography}

\end{document}
